% !TEX program = xelatex
\documentclass[12pt,a4paper]{article}

\usepackage[UTF8]{ctex}
\usepackage{geometry}
\geometry{left=2.5cm, right=2.5cm, top=2.5cm, bottom=2.5cm}
\usepackage{hyperref}
\hypersetup{colorlinks=true, linkcolor=blue, citecolor=blue, urlcolor=blue}
\usepackage{enumitem}
\usepackage{booktabs}
\usepackage{array}
\usepackage{float}
\usepackage{titlesec}

\titleformat{\section}{\Large\bfseries}{\thesection}{1em}{}
\titleformat{\subsection}{\large\bfseries}{\thesubsection}{1em}{}
\titleformat{\subsubsection}{\normalsize\bfseries}{\thesubsubsection}{1em}{}

\begin{document}

\begin{center}
    {\LARGE\bfseries 煤矿巷道锚杆支护系统的预紧力矩优化设计}\\[1em]
    {\large AI工具使用说明}
\end{center}

\vspace{1em}

本作品在撰写过程中使用了AI工具辅助完成部分数据分析与建模工作。以下详细说明所用AI工具的名称、版本、使用目的、关键交互记录及人工修改情况。

% ============================================================
\section{所用AI工具信息}

\begin{enumerate}
    \item \textbf{DeepSeek网页端}（\url{https://chat.deepseek.com}）
    \begin{itemize}
        \item 模型版本：DeepSeek-V3 / DeepSeek-R1
        \item 使用时间：2026年7月18日--7月21日
        \item 用途：统计检验方法咨询、岩石力学参数查询、数据处理算法验证
    \end{itemize}

    \item \textbf{集成开发环境（IDE）内置AI辅助}
    \begin{itemize}
        \item 工具：Visual Studio Code + AI代码补全插件
        \item 用途：Python数据分析脚本编写、LaTeX公式排版辅助
        \item 说明：IDE内置AI辅助功能通过代码补全和语法建议提升编码效率，不涉及独立对话记录导出。
    \end{itemize}
\end{enumerate}

% ============================================================
\section{具体使用目的与环节}

\begin{table}[H]
\centering
\caption{AI工具使用情况总览}
\label{tab:ai-usage}
\renewcommand{\arraystretch}{1.3}
\begin{tabular}{|c|p{5.5cm}|p{5.5cm}|c|}
\hline
\textbf{序号} & \textbf{使用目的} & \textbf{对应环节} & \textbf{工具} \\
\hline
1 & Bootstrap法重抽样置信区间估计 & 问题1临界预紧力矩的统计检验 & DeepSeek网页端 \\
\hline
2 & 岩石力学参数（弹性模量、普氏系数等）文献查询 & 问题4参数标定与文献数据整理 & DeepSeek网页端 \\
\hline
3 & Chow检验（结构断点检验）原理咨询与方案评估 & 问题1.1分段回归临界点检验方法选型 & DeepSeek网页端 \\
\hline
4 & Python数据处理脚本编写（分段回归、敏感性分析等） & 全文数值计算 & IDE内置AI \\
\hline
5 & LaTeX公式与图表排版 & 论文撰写阶段 & IDE内置AI \\
\hline
\end{tabular}
\end{table}

% ============================================================
\section{关键交互记录}

\subsection{Bootstrap法置信区间估计}

对话链接：\url{https://chat.deepseek.com/share/cbztzpcwa78v7rj4se}

\begin{itemize}
    \item \textbf{主要提问内容}：\\
    请帮我理解Bootstrap法的基本原理，并给出用Python实现对回归残差进行Bootstrap重抽样、估计参数置信区间的代码框架。
    
    \item \textbf{AI回复要点}：
    \begin{enumerate}
        \item 解释了Bootstrap法通过有放回重抽样模拟总体分布的统计思想；
        \item 给出了基于残差重抽样的Bootstrap回归置信区间估计的Python实现代码；
        \item 说明了重抽样次数（如$B=1000$）对结果稳定性的影响。
    \end{enumerate}
    
    \item \textbf{采纳与修改情况}：\\
    采纳了Bootstrap法的基本框架与代码结构，根据本课题数据特点进行了以下修改：
    \begin{itemize}
        \item 将通用回归模型替换为过原点线性回归$T = a \cdot P$的特定形式；
        \item 调整了重抽样次数与置信水平参数；
        \item 将输出结果整合到本文分段回归的临界点检验流程中。
    \end{itemize}
\end{itemize}

\subsection{岩石力学参数查询}

对话链接：\url{https://chat.deepseek.com/share/4fdrjwepoeweoo8ty6}

\begin{itemize}
    \item \textbf{主要提问内容}：\\
    请提供煤矿巷道常见围岩（砂岩、页岩、煤体等）的弹性模量$E$和普氏系数$f$的典型取值范围及相关文献出处。
    
    \item \textbf{AI回复要点}：
    \begin{enumerate}
        \item 列举了砂岩、页岩、泥岩、煤体等岩性的弹性模量与普氏系数典型范围；
        \item 提供了部分参考文献的基本信息；
        \item 说明了参数随含水状态、风化程度等因素的变化规律。
    \end{enumerate}
    
    \item \textbf{采纳与修改情况}：\\
    采纳了AI提供的参数范围作为初步参考，后续通过查阅矿业工程领域权威文献进行了交叉验证与修正，确保数据来源可靠。最终用于问题4幂律拟合$E(f) = \beta \cdot (10f)^\alpha$的数据集经过人工筛选与核实。
\end{itemize}

\subsection{Chow检验（结构断点检验）}

对话链接：\url{https://chat.deepseek.com/share/ayiw8vzp39u6jzbmb3}

\begin{itemize}
    \item \textbf{主要提问内容}：\\
    请解释Chow检验的统计原理，以及如何用Python实现对分段线性回归的结构断点检验，检验临界预紧力矩处回归系数是否发生显著变化。
    
    \item \textbf{AI回复要点}：
    \begin{enumerate}
        \item 详细解释了Chow检验的原假设（$H_0$：断点前后回归系数无显著差异）与$F$统计量构造方法；
        \item 给出了基于OLS回归的Chow检验Python实现代码；
        \item 说明了检验结果的判读方法（$p$值与显著性水平的关系）。
    \end{enumerate}
    
    \item \textbf{讨论与决策过程}：\\
    团队在与AI讨论Chow检验原理后，进一步在队伍内部进行了深入研讨，最终\textbf{决定不在论文中采用该方法}，理由如下：
    \begin{itemize}
        \item Chow检验要求断点前后回归模型形式一致（仅系数不同），而本课题中临界点两侧的物理机制存在差异（稳定段与非线性段），模型假设不完全满足；
        \item 本课题已采用瞬时$K$滑动方差收敛法作为主检验方法，辅以Bootstrap置信区间估计，两种方法交叉验证已足够稳健；
        \item 经团队讨论认为，保留物理意义更明确的滑动方差法比引入统计假设受限的Chow检验更为合理。
    \end{itemize}
    该决策体现了团队独立思考与科学判断能力：在充分了解AI提供的方法后，基于问题本身的物理背景做出了合理取舍。
\end{itemize}

% ============================================================
\section{人工修改与质量保证}

\begin{enumerate}
    \item \textbf{数据层面}：所有通过AI获取的岩石力学参数均经过原始文献查阅与交叉验证，确保数据准确性。
    
    \item \textbf{算法层面}：AI提供的Bootstrap法代码仅作为方法学参考，实际实现中根据本课题的过原点回归模型进行了针对性修改，并通过多组数据验证了算法正确性。
    
    \item \textbf{论文层面}：全文由人工独立撰写，AI工具仅辅助完成技术细节的理解与代码实现，未直接生成论文文本内容。所有公式推导、结果分析与结论均经过作者独立验证。对于AI提供的Chow检验方法，团队在充分理解其原理后，经内部讨论决定不纳入最终论文，体现了独立判断能力。
    
    \item \textbf{交叉验证}：临界预紧力矩$T_c$的判定采用了Bootstrap置信区间与瞬时$K$滑动方差收敛法两种独立方法交叉验证，确保结论稳健可靠。
\end{enumerate}

\vspace{1em}
\noindent\textbf{声明}：本团队对论文中所有内容的真实性、准确性和原创性负责。AI工具的使用仅限于辅助理解统计方法原理、查询参数参考范围与辅助编写计算代码，核心建模思路、数据分析过程与论文撰写均由团队成员独立完成。

\end{document}
